\documentclass[11pt]{article}
\usepackage[a4paper,margin=25mm]{geometry}
\usepackage{amsmath,amssymb,amsthm,booktabs,longtable,hyperref,xurl,microtype}
\setlength{\emergencystretch}{3em}
\hypersetup{hidelinks}
\newtheorem{theorem}{Theorem}
\newtheorem{lemma}[theorem]{Lemma}
\newtheorem{proposition}[theorem]{Proposition}
\newtheorem{corollary}[theorem]{Corollary}
\newcommand{\Q}{\mathbb Q}\newcommand{\Z}{\mathbb Z}
\title{Common physical points in a combined quadratic system\\
Quartic checkpoints in 1601 and 1837 variables}
\author{\(\mathbb Z\mathbb Q\) Diophantine Project}
\date{23 November 2026}
\begin{document}
\maketitle
\begin{abstract}
We give complete quadratic systems with 1601 rational coordinates and
1599 rows, or 1837 coordinates and 1811 rows. Their sums of squares are
quartics. The construction identifies the actual common elliptic point
functions in twelve integrality blocks: 54 stored coordinates in the
general 163-coordinate interface, or 50 in its specified 179-coordinate
counterpart. A full physical-function dictionary distinguishes those
functions from variables carrying input-dependent or contact-dependent
scales. An additional common input numerator is justified by the
all-sign order coprimality of the retained finite prime supplier and a
high-precision Chinese remainder construction.

The finite counts, maps, and duplicate equations are verified exactly.
The undecidability consequences remain conditional on all retained
arithmetic interfaces and on the separately unresolved universality
assertion \(U_{\mathrm{em}}\) for the exact emended integer source.
No unconditional undecidability bound or global minimum is asserted.
\end{abstract}

\section{Scope and predecessor systems}

The starting rational systems are the complete reciprocal-reference
checkpoints with 2085 and 2277 coordinates \cite{reciprocal}. Their
simultaneous witness construction was proved in the shared-reference
checkpoint \cite{shared}, and their outer integer system and its
reconstruction are recorded in \cite{outer}. The present paper changes
which existing integrality witnesses are shared. The source coordinate
maps and full resulting coefficient tables are part of the deliverable.

The source-universality obligation \(U_{\mathrm{em}}\) means universality
of the precise emended positive-integer system used in \cite{outer}.
The consulted 1980 display prints Greek \(\alpha\) in a coefficient
where the emendation uses Latin \(a\). The later survey supports the
Latin coefficient but has a separate impossible Pell-row equality.
Consequently no fully matching published display has been verified,
and that obligation remains open. All algebraic equivalences and
coordinate counts below are independent of \(U_{\mathrm{em}}\).
It is invoked only in the final undecidability implication.

The arithmetic interfaces retained from the selected integer-defining
source include its elliptic-group and point-chart interpretations,
height and contact bounds, finite-family prime supplier, compatible
finite-place and compact-real recurrence statements, and the stated
native-pole realization. The common-index step uses the supplier's
all-sign reduction-order coprimality, audited in \cite{crt}.
These general routes add no ordinary rational \(abc\) hypothesis;
the statement that they retain other arithmetic premises is essential.
The complete proofs of those interfaces are not supplied by a finite
coefficient verification.

\section{Which physical data may be shared}

Write \(K=\mathbb Q(\sqrt2)\), with nontrivial automorphism denoted by a
bar.  All field quantities below are represented by two rational
coordinates.  Fixed factors of two in the source's imaginary-component
storage convention are retained throughout.  They have no dependence on
the input or on a contact parameter.

The simultaneous completeness proof in the shared-reference checkpoint
chooses one reference point \(H\), one positive scalar label \(a\), and
one pair of Lucas labels \(t,b\), for any fixed finite list of integer
inputs.  Consequently the difference points \(Q_2,Q_3,Q_4\) are common.
The numerator pairs used for these three roles can also be chosen common:
they are the source multiples of that same \(H\), with the same labels.
The input-dependent difference \(Q_{1,i}\) and the contact parameter of
block \(i\) remain independent.  No new common-point existence assertion
is being inferred merely from equality of row names.

\subsection{Reference, point, and Lucas functions}
Put
\[
 f=\frac1{x(H)},\qquad S=\frac{y(H)}{x(H)},\qquad
 T_H=\frac{y(H-P_0)}{x(H-P_0)}.
\]
The common reference functions are
\[
      f,\quad S,\quad f(H-P_0)^2,\quad f^2,\quad T_H.
\]
Their ten rational coordinates were already shared in the previous
checkpoint.

For a finite numerator pair in role \(j\in\{1,3,4\}\), write
\[
 s_j=\frac{y(U_j)}{x(U_j)},\quad
 t_j=\frac{y(V_j)}{x(V_j)},\quad
 F_j=\frac1{x(U_j)},\quad G_j=\frac1{x(V_j)}.
\]
Let \(\eta_j,d_j\) be the other two fields in the finite-pair chart,
and set \(p_j=\eta_j+s_j+t_j\).  They satisfy
\begin{align}
 d_j&=p_jF_j=\eta_jG_j,\label{eq:common-finite-pair-1}\\
 D(F_j+G_j)&=(s_j-t_j)\eta_j-t_j^2,\label{eq:common-finite-pair-2}\\
 1&=d_j(s_j-t_j)-DF_jG_j.\label{eq:common-finite-pair-3}
\end{align}
Here \(D=11-6\sqrt2\).  These are precisely the finite-pair equations
proved in the physical chart manuscript \cite{physical}.  On the
completeness witnesses the six fields are rational functions of the
chosen pair, with the tangent case interpreted by the finite chart.
For roles 3 and 4 every one of them is common.  We share only the fields
whose existing stored coordinates already have the required scale.

For role 2 use the normalized inverse-Vieta fields
\((u,v,w,\rho)\) of \cite{scalar}, and put
\[
 Z=\frac{2z(Q_2)}{z(H)},\qquad
 (U,V,W)=Z(u,v,w),\qquad r_2^*=Z^2\rho^2.
\]
The retained primary ratio is \(A=L_2\).  The other two point functions
needed below are
\[
                F_{U_2}=f(U_2)^2,\qquad F_{V_2}=f(V_2)^2.
\]
All of \(A,Z,U,V,W,r_2^*,F_{U_2},F_{V_2}\) depend only on the
common role-2 numerator pair and on \(H\).  In particular, none depends
on the rational scalar normalizer \(R_i\).  This distinction matters:
\(P_i^*=R_i x(Q_2)\) does depend on \(R_i\) and is not shared.

The two common cross-errors are
\[
 E_3=S\overline{t_3}-\overline{T_H}s_3,
 \qquad E_4=S\overline{t_4}-\overline{T_H}s_4.
\]
With \(w_4=s_4\), define the common Lucas auxiliary fields
\begin{equation}\label{eq:common-lucas-fields}
 \mathcal C=(A-1)(w_4+4S),\qquad
 \mathcal P s_3=2Sw_4,\qquad
 \mathcal Q=(A+1)\mathcal P.
\end{equation}
On the completeness witnesses the finite exclusions ensure \(s_3\ne0\),
so \(\mathcal P\) is the same in all blocks.  No division by \(s_3\)
is added to the equations.  Soundness will follow by restoring each
original block separately, and therefore includes every allowed
zero-ordinate boundary of the original system.

\subsection{The exact physical-coordinate dictionary}

Use \(X_j\) for a physical source coordinate and \(Y_j\) for its stored
coordinate before the final renumbering.  The source maps have the form
\begin{equation}\label{eq:physical-laurent-map}
                 X_j=c_i^{\omega_j}R_i^{\rho_j}
                       C^{\kappa_j}Y_j,
\end{equation}
where \(c_i\ne0\) is block \(i\)'s contact gauge.  The parameter \(C\)
is a common fixed coefficient, and is 1 in the fixed interface.  We do
not identify any of the \(c_i\) or \(R_i\).

The following table gives every directly shared physical source pair;
the next section adds a common input numerator.
A dash means that the current stored scale is not shared in that
interface.  Each listed pair contributes two rational coordinates.
The \(c_i\)-normalization column gives the multiplier applied to its physical
source pair to obtain the displayed function; a fixed
factor \(C\), where present, is also shown.  All listed pairs have
\(R_i\)-weight zero.

\begin{center}\small
\begin{longtable}{@{}p{34mm}p{41mm}c c c@{}}
\toprule
Physical function & Source coordinate pairs & Normalization & 163 & 179\\
\midrule\endfirsthead
\toprule
Physical function & Source coordinate pairs & Normalization & 163 & 179\\
\midrule\endhead
Reference functions & \(\begin{gathered}15\!:\!16,17\!:\!18,19\!:\!20\\64\!:\!65,150\!:\!151\end{gathered}\)
 & \(c_i^{-1}\) & 10 & 10\\
\(A\) & \(33\!:\!34\) & \(C/c_i\) & 2 & 2\\
\(Z\) & \(102\!:\!103\) & 1 & 2 & 2\\
\(U,V,W\) & \(182\!:\!183,184\!:\!185,186\!:\!187\) & 1 & 6 & 6\\
\(F_{V_2}\) & \(23\!:\!24\) & \(c_i^{-1}\) & 2 & --\\
\(F_{U_2}\) & \(54\!:\!55\) & \(C/c_i\) & 2 & --\\
\(r_2^*\) & \(194\!:\!195\) & 1 & 2 & --\\
\(s_3,t_3,\eta_3\) & \(213\!:\!214,215\!:\!216,221\!:\!222\) & 1 & 6 & 6\\
\(F_3\) & \(217\!:\!218\) & 1 & -- & 2\\
\(s_4,t_4,\eta_4\) & \(225\!:\!226,227\!:\!228,233\!:\!234\) & 1 & 6 & 6\\
\(d_4\) & \(235\!:\!236\) & 1 & 2 & --\\
\(G_4\) & \(231\!:\!232\) & 1 & -- & 2\\
\(E_3,E_4\) & \(39\!:\!40,43\!:\!44\) & \(c_i^{-1}\) & 4 & 4\\
\(\mathcal P,\mathcal C,\mathcal Q\) & \(41\!:\!42,46\!:\!47,48\!:\!49\)
 & \(c_i^{-1}\) & 6 & 6\\
\midrule
Total without common input numerator &&&50&46\\
\bottomrule
\end{longtable}
\end{center}

For example, the physical primary ratio is
\[
 A=\frac C{c_i}(X_{33}+\sqrt2X_{34}).
\]
The general map has \(\omega_{33}=\omega_{34}=1\),
\(\rho_{33}=\rho_{34}=0\), and
\(\kappa_{33}=\kappa_{34}=-1\), so
\(A=Y_{33}+\sqrt2Y_{34}\).  Thus it is the actual common ratio,
not a ratio multiplied by the block's gauge, that is identified.
Similarly
\[
 F_{U_2}=\frac C{c_i}(X_{54}+2\sqrt2X_{55})
             =Y_{54}+2\sqrt2Y_{55}
\]
in the 163-coordinate interface.  For \(Z,U,V,W\), every relevant
\(\omega_j\) and \(\rho_j\) is zero.  The same is true of all listed
finite-pair fields.  For each reference, error, and Lucas field its
factor \(c_i^{-1}\) cancels its stored \(c_i\)-weight one.
The displayed weights and parameter powers are checked directly
against the pinned \texttt{maps.json} files.

One must not treat every physical common function as already shareable.
For instance the 163-coordinate chart stores the components of \(\rho Z\)
with nonzero \(R_i\)-weight, and some fields in the Lucas square circuit
retain contact-gauge factors.  These coordinates are kept independent.
Likewise the 179-coordinate storage of \(F_{U_2},F_{V_2},r_2^*\) has a
different contact scale, so those pairs are omitted there.

\begin{proposition}[Direct common-point sharing]\label{prop:common-direct}
Under the same retained arithmetic interfaces as the shared-reference
construction, one may identify the 50 or 46 coordinates in the table
in any finite list of blocks.  Every solution of the identified system
has integral distinguished inputs, and every tuple of integers has a
rational witness of the identified system.
\end{proposition}
\begin{proof}
For soundness, give each original block its own copy of each identified
value.  Every original residual is recovered; no equation has been
weakened.  The original defining theorem makes each distinguished
input integral.

For completeness, apply the simultaneous scalar preparation and contact
separation from the shared-reference proof.  It supplies one final \(H\)
and common labels for roles 2, 3, and 4, and therefore the common numerator
pairs used above.  The required finite chart exclusions are all exclusions
of finitely many points, and were already allowed when the common
reference neighborhoods were chosen.  Fill all chart coordinates from
these actual points.  The displayed physical dictionary shows that every
coordinate being identified then has the same rational value in all
blocks.  The normalizers, contact gauges, contact slopes, and native-pole
witnesses retain their independent values.  Consequently every original
block remains satisfied, including all its local preparations, and the
identified system has the required witness.
\end{proof}

\section{Sharing the input numerator across disjoint contact families}

There is a further four-coordinate saving per additional block.  It is
not justified by the preceding common scalar labels, because the input
root labels
\[
                      \beta_i=51x_i^2+2504
\]
usually differ.  Instead we use the all-sign coprimality part of the
retained prime supplier and choose a common numerator after the contact
families have been separated.

\begin{lemma}[A common input numerator]\label{lem:common-input-numerator}
In the simultaneous completeness construction one may choose
\[
               U_{1,i}=U\qquad\text{for every }i,
\]
while retaining the common reference \(H\), the common roles 2, 3, and 4,
and the independent differences \(Q_{1,i}=[\beta_i]P_0\).
The alternate points are \(V_{1,i}=U-Q_{1,i}\), and are not identified.
\end{lemma}

\begin{proof}
We record both of the coprimality assertions that are needed.  Let \(d\)
run through the finite union of translated contact positions, and let
\(p_d\) be its fresh prime.  After the prescribed Gaussian choice, set
\[
                 u_d=N(\pi_d-1)/4,
                 \qquad v_d=N(\pi_d+1)/4.
\]
The source residue conditions give \(v_2(u_d)=4\) and \(v_2(v_d)=0\).
The all-sign order assertion of the retained supplier states that at
different positions every common divisor of the relevant two reduction
orders is a power of two.  Together with its same-position assertion
this implies
\begin{equation}\label{eq:common-input-orders}
 \gcd(u_d,v_e)=1\quad\text{for all }d,e,
 \qquad \gcd(v_d,v_e)=1\quad\text{if }d\ne e.
\end{equation}
The separate short-interval argument gives
\[
              \gcd\!\left(p_d,\prod_eu_ev_e\right)=1.
\]
The latter assertion permits arbitrary additional prime-power precision;
it does not replace \eqref{eq:common-input-orders}.  The supplier and
this precise distinction are audited in the project's common-index
CRT argument \cite{crt}.

Before choosing the contact primes, retain the original scalar
preparation for \(H\) and choose a common numerator neighborhood on
which the finitely many pairs
\((U,U-Q_{1,i})\) have valid finite charts.  Only finitely many points
are excluded: identity points, forbidden abscissas, zero ordinates when
needed, and any denominator of a retained rational point function.
An infinite cyclic orbit therefore supplies a point avoiding their
union.  Choose its real and prepared-place neighborhoods sufficiently
small that all the relevant raw functions are bounded.  Include the
associated progression modulus \(M\), and the compact-component order,
in finite avoidance.  For precision about the source supplier, its
odd avoidance modulus excludes the prime seven, but seven is
already avoided automatically: its prescribed residue
\(\pi_d\equiv1+3\mathrm i\pmod7\) gives
\(N(\pi_d-1)\equiv2\pmod7\) and
\(N(\pi_d+1)\equiv6\pmod7\).  All other fixed odd primes
are handled by the ordinary avoidance condition.  We may take \(64\mid M\) and choose its allowed
index residue \(j_0\) divisible by 64.  The odd parts of the selected
orders are then prime to \(M\).

Separate the contact lists by the fixed translations already used in
the shared-reference proof.  In particular, each contact belongs to
exactly one input block.  The aggregate shifts are unchanged: they
depend on the differences \(Q_{1,i},Q_2,Q_3,Q_4\), which were fixed
before the new numerator was chosen.  Apply the retained supplier to
this union, including its all-sign order conditions.

Complete the original choice \(H=[h]P_0\), with its annihilation
congruences and its synchronized primary and opposite formal depths.
For a contact \(d\) of block \(i\), impose on one integer \(j\)
\begin{equation}\label{eq:common-input-crt}
 \begin{aligned}
 j&\equiv j_0\pmod M,\\
 j&\equiv0\pmod{u_d},\qquad
 j\equiv\beta_i\pmod{v_d},\\
 j&\equiv\beta_i h\pmod{p_d^{L_d}}.
 \end{aligned}
\end{equation}
Here the positive integers \(L_d\) will be chosen sufficiently large.
The primary-order congruences all have residue zero.  The opposite
orders are pairwise coprime and prime to all primary orders by
\eqref{eq:common-input-orders}.  Their only possible interaction with
\(M\) has already been excluded; the primary 2-parts are compatible
because \(64\mid j_0\).  Finally the selected prime powers are coprime
to these orders, to \(M\), and to one another.  The generalized Chinese
remainder theorem therefore supplies one progression for \(j\).
The retained compact-component recurrence gives positive \(j\) in this
progression with \([j]P_0\) in the previously chosen real neighborhood.
All prepared finite-place neighborhoods are preserved by its modulus
\(M\).

Put \(U=[j]P_0\).  At a contact of block \(i\), the primary reduction
of \(U\) is the identity and the opposite reduction of
\(U-Q_{1,i}\) is the identity.  Moreover \(j-\beta_i h\) is divisible
by the relevant annihilator at both slots and by \(p_d^{L_d}\).
Multiplication by increasing powers of \(p_d\) in the formal group
therefore makes
\[
                     U-[\beta_i]H
\]
arbitrarily close to the identity at both slots.  Consequently
\((U,U-Q_{1,i})\) is arbitrarily close to the original pair
\(([\beta_i]H,[\beta_i](H-P_0))\) at that contact.

This last statement preserves every inequality actually used in the
source proof.  With the reference point fixed, the finitely many raw
chart functions are continuous at its chosen finite local points.
Their required lower valuation bounds are open conditions; this
also holds when an original error is exactly zero.  Choose \(L_d\)
large enough to preserve the reciprocal depths, the input-numerator
depth, the paired-comparison depth, and the resulting strict factor
inequalities.  Equivalently, in formal logarithms the primary and
opposite quotients differ from \(\beta_i\) by
\[
       \frac{j-\beta_i h}{h},\qquad
       \frac{j-\beta_i h}{h-1},
\]
whose valuations can exceed any fixed required precision.  The
relative fourth-order estimates and the synchronized reference depths
are thus exactly those used in the original construction.  The
comparison gates still give native-pole input valuation one at their
selected contacts.

At the fixed finite and real places no common ratio is required to be
close to the distinct values \(\beta_i\).  The physical tests are
bounded on the common neighborhoods just chosen.  The source contact
expansion
\[
       q(\lambda+u)=\kappa\lambda^{-2}+O(\lambda^{-3})
\]
holds simultaneously for all finitely many fixed shifts.  All test
powers in the physical comparison factors have positive contact
powers, hence the factors tend to one as the common slope grows at
these places.  The source calculation gives \(W_i\to1\), gate error
\(D_i\to0\), and the required actual native-pole limits and square
classes.  The scalar preparations do not involve \(U\) and are
unchanged.  The unselected-prime parity argument is unchanged as well.
Every actual native pole can therefore be filled with its own
realization witnesses.  This proves the lemma.
\end{proof}

\begin{corollary}\label{cor:four-more}
In addition to the preceding dictionary, one may identify physical
source coordinates \(201,202,205,206\), representing \(s_1,F_1\), in
both general interfaces.  There are therefore 54 common stored
coordinates in the 163-coordinate interface and 50 in the
179-coordinate interface.
\end{corollary}
\begin{proof}
Both source pairs have \(c_i\)-weight zero, \(R_i\)-weight zero, and
parameter-weight zero in both published maps.  They are exactly
\(s_1=y(U)/x(U)\) and \(F_1=1/x(U)\), common by the lemma.
Their alternate partners \(t_{1,i},G_{1,i}\) and the corresponding
finite-pair \(\eta_{1,i},d_{1,i}\) remain independent.  Soundness still
uses each recovered original block separately.
\end{proof}

\section{Complete combined systems and counts}

The optimized outer system has 251 coordinates and 219 quadratic rows.
It uses twelve integer blocks, including the reciprocal block of the
preceding checkpoint.  That reciprocal change modifies only local
registers \(37,77,78\), none of which is shared here.  The three earlier
input-square identifications also change none of the present shared
registers.

For an interface with \(m\) coordinates and \(q\) rows, let \(s\) be
the number of common coordinates and \(r\) the number of rows supported
entirely on these coordinates that become literal duplicates.  With
\(k\) integer inputs, the combined counts are
\begin{equation}\label{eq:common-count}
 N=251+k(m-1)-s(k-1),\qquad
 Q=219+kq-r(k-1).
\end{equation}
The constant 251 is specific to this outer circuit; the same argument
replaces it by any other compatible outer coordinate count.

For the 163-coordinate interface, \(s=54\) and \(r=24\).  Eight
reference rows were already repeated.  The sixteen additional repeated
rows are the two components of each of
\[
 E_3,\ E_4,\ \mathcal C,\ \mathcal Q,\ \mathcal P s_3-2Ss_4,
 \ F_{U_2}r_2^*-U^2,\ F_{V_2}r_2^*-V^2,
 \ A(U+M)-2SZ,
\]
where \(M=-2S-W\) in the inverse-Vieta chart.
For the 179-coordinate interface \(s=50\) and \(r=20\); the four
rows involving \(F_{U_2},F_{V_2},r_2^*\) retain their private storage
and are not duplicates.  Sharing \(s_1,F_1\) adds no duplicate row,
because the input pair still contains independent alternate fields.

Putting \(k=12\) in \eqref{eq:common-count} gives
\[
\begin{array}{c|rrrr|rr}
 m&q&s&r&k&N&Q\\\hline
 163&137&54&24&12&1601&1599\\
 179&151&50&20&12&1837&1811.
\end{array}
\]
If the additional common-input lemma is not used, the direct dictionary
alone gives \(1645\) and \(1881\) coordinates with the same respective
row counts.  The earlier independent-block and shared-reference
checkpoints remain separately reproducible.

\begin{theorem}[Conditional quartic transfer]\label{thm:common-transfer}
Retain all arithmetic interfaces of the respective integer-defining
construction, including the finite-family supplier and the all-sign
order assertion used above.  The complete combined systems just
constructed have 1601 or 1837 rational coordinates and 1599 or 1811
quadratic residuals, respectively.  Their sums of squares have degree
exactly four and the same outer input projection as the emended
positive-integer source system on the retained cofinite parameter set.

If, additionally, the source-universality assertion
\(U_{\mathrm{em}}\) holds for that precise emended system, these
families give the corresponding undecidability reductions.  The
universality assertion is not proved by this construction, and no
unconditional undecidability bound or minimality assertion is made.
\end{theorem}
\begin{proof}
For a rational solution, restore all duplicated coordinates and rows.
The two rational coordinate changes and every unchanged native block
then recover the original independent certificates.  Their soundness
makes the required source inputs integral; the already proved outer
reconstruction supplies a positive-integer solution of the emended
source.  Conversely, start with a positive-integer source witness,
apply the outer sign and reciprocal transformations, and use
Proposition~\ref{prop:common-direct} and
Lemma~\ref{lem:common-input-numerator} to fill the twelve blocks.
All shared coordinates agree by their physical dictionary.  This
proves equality of projections without using \(U_{\mathrm{em}}\).
The latter assertion is needed only for the undecidability consequence.

Every residual is still quadratic: identification replaces variables
by variables, and only literal duplicate rows are omitted.  At least
one retained row, for example the reciprocal input-unit row, has a
nonzero quadratic part.  Over \(\mathbb Q\) the sum of squares has
zero set exactly the common zero set, and its degree is exactly four
because a sum of real polynomial squares cannot cancel its nonzero
leading quadratic squares.
\end{proof}

\paragraph{Exact verification.}
The generator starts from the complete reciprocal-reference tables,
reads the literal source-to-stored maps, and checks the contact and
normalizer weights of every identified coordinate.  It materializes
every resulting residual, records every removed duplicate and its
retained identical row, and checks the full coordinate support and
quadratic degree.  The separate verifier replays the identification on
the complete input table, checks every kept and deleted row, and checks
the exact final bytes and hashes.  These finite checks do not verify
\(U_{\mathrm{em}}\), the external arithmetic interfaces, or a global
minimum.

\begin{thebibliography}{9}
\bibitem{outer} ZQ Diophantine Project,
\emph{The emended Jones outer system and its quartic transfer},
\nolinkurl{undecidability/papers/jones1980_quartic_transfer.tex},
23 November 2026. See especially the source-reconciliation section.
\bibitem{shared} ZQ Diophantine Project,
\emph{A shared-reference quartic checkpoint with an explicit
source-universality obligation},
\nolinkurl{undecidability/papers/shared_quartic_checkpoint.tex},
23 November 2026.
\bibitem{reciprocal} ZQ Diophantine Project,
\emph{A reciprocal chart saving one coordinate in the shared-reference
quartic construction},
\nolinkurl{undecidability/papers/reciprocal_reference_checkpoint.tex},
23 November 2026.
\bibitem{crt} ZQ Diophantine Project,
\emph{A common numerator point for simultaneous integer certificates},
\nolinkurl{undecidability/papers/common_numerator_transfers.tex},
23 November 2026. See the all-sign order assertion and its
fixed-modulus compatibility, including the automatic avoidance of seven.
\bibitem{physical} ZQ Diophantine Project, based on the retained source
notes of Gabe Conant and the project's subsequent chart reductions,
\emph{A quartic in 193 rational coordinates},
\nolinkurl{papers/quartic/quartic193_20261117_v74.tex},
17 November 2026; complete source-pair definitions in
\nolinkurl{computation/quartic/seven_units193_v74/build.py}.
\bibitem{scalar} ZQ Diophantine Project,
\emph{Scalar and point reductions to 163 and 179 coordinates},
\nolinkurl{papers/quartic/scalar_point163_179_20261120_v82.tex},
20 November 2026; independently reconstructed physical equations in
\nolinkurl{computation/variable_reductions_v82/general/check.py}.
\end{thebibliography}
\end{document}
